\chapter{EXPERIMENTAL RESULTS AND ANALYSIS}\label{chap:results}

This chapter presents the results of the full experimental evaluation. Section~\ref{sec:synth_results} covers the 2,160-row synthetic benchmark; Section~\ref{sec:real_results} covers the 300-row real-data validation on RetailRocket and RecSys~2015.

\section{Synthetic Experiment Results}\label{sec:synth_results}

\subsection{Speedup vs.\ Pruning Threshold $\tau$}

Table~\ref{table:speedup_tau} summarises median speedups aggregated across all graph types, sizes, degrees, and distributions.

\begin{table}[h!]
\centering
\renewcommand{\arraystretch}{1.3}
\begin{tabular}{@{}rrrrrr@{}}
\toprule
$\tau$ & \textbf{Edge Speedup} & \textbf{Wall Speedup} & \textbf{Baseline (ms)} & \textbf{Pruned (ms)} & \textbf{Path Found (\%)} \\
\midrule
0.001 &   8.8$\times$ &   3.6$\times$ & 1.67 & 0.496 & 28.6 \\
0.01  &  82.1$\times$ &  30.4$\times$ & 1.67 & 0.055 & 5.8  \\
0.05  & 374.4$\times$ & 146.8$\times$ & 1.67 & 0.011 & 4.7  \\
0.1   & 790.8$\times$ & 280.9$\times$ & 1.67 & 0.006 & 4.4  \\
0.5   & 2,124$\times$ & 881.3$\times$ & 1.67 & 0.002 & 3.9  \\
\bottomrule
\end{tabular}
\caption[Speedup by pruning threshold $\tau$]{Median speedups by $\tau$ across all synthetic configurations. ``Edge Speedup'' = baseline edges relaxed / pruned edges relaxed. ``Wall Speedup'' = baseline ms / pruned ms. Optimality gap $\Delta = 0.0000\%$ for all path-found rows.}
\label{table:speedup_tau}
\end{table}

Speedups scale monotonically with $\tau$, ranging from $3.6\times$ wall-clock at $\tau=0.001$ to $881\times$ at $\tau=0.5$. Edge speedups consistently exceed wall-clock speedups (e.g., $2{,}124\times$ vs.\ $881\times$ at $\tau=0.5$) because Python per-call overhead becomes a larger fraction of the substantially reduced pruned runtime. Across all 171 path-found rows among the 1,800 pruned synthetic runs ($\tau > 0$), the optimality gap remained exactly $\Delta = 0.0000\%$, confirming the all-or-nothing property (Theorem~\ref{thm:all_or_nothing}, Chapter~\ref{chap:correctness}). Path-found rates decrease with $\tau$, reflecting the reachability--speed trade-off inherent in the pruning approach.

\begin{figure}[ht]
\centering
\includegraphics[width=\textwidth]{plot1_speedup_vs_tau.png}
\caption[Speedup vs.\ pruning threshold $\tau$]{Wall-clock speedup as a function of $\tau$ for Erd\H{o}s--R\'{e}nyi (left) and layered (right) graphs.}
\label{fig:speedup_vs_tau}
\end{figure}

Figure~\ref{fig:speedup_vs_tau} plots the wall-clock speedup ratio (baseline time divided by pruned time) against the pruning threshold $\tau$ on a logarithmic $x$-axis spanning $\tau \in [0.001, 0.5]$. Each line represents a unique combination of graph size $|V|$ and edge-weight distribution. Speedup increases monotonically with $\tau$ across all configurations, confirming that more aggressive pruning yields proportionally faster search. Erd\H{o}s--R\'{e}nyi graphs consistently achieve higher speedups than layered graphs because their uniformly random connectivity exposes a larger fraction of low-probability edges to the threshold. In contrast, the layered generator's built-in forward bias toward the target means that many edges carry relatively high probability and therefore survive pruning. Larger graphs (higher $|V|$) produce steeper speedup curves, reflecting the greater absolute fraction of the search space that the threshold eliminates as the graph grows.

\subsection{Speedup vs.\ Graph Size $|V|$}

Table~\ref{table:speedup_size} shows how speedup scales with graph size at fixed $\tau = 0.01$.

\begin{table}[h!]
\centering
\renewcommand{\arraystretch}{1.3}
\begin{tabular}{@{}rrrrrr@{}}
\toprule
$|V|$ & \textbf{Edge Speedup} & \textbf{Wall Speedup} & \textbf{Baseline (ms)} & \textbf{Pruned (ms)} & \textbf{Peak Mem (KB)} \\
\midrule
1,000  &  30.6$\times$ & 11.0$\times$ & 0.58 & 0.050 & 83.9 $\to$ 4.4 \\
5,000  & 104.9$\times$ & 55.8$\times$ & 2.69 & 0.056 & 349.6 $\to$ 4.6 \\
10,000 & 173.9$\times$ & 89.1$\times$ & 5.64 & 0.057 & 662.3 $\to$ 4.5 \\
\bottomrule
\end{tabular}
\caption[Speedup scaling by graph size]{Median speedups at $\tau = 0.01$ by graph size. Baseline time grows linearly with $|V|$; pruned time remains nearly constant. Memory columns show baseline $\to$ pruned peak consumption.}
\label{table:speedup_size}
\end{table}

Baseline execution time scales linearly with $|V|$ ($0.58 \to 2.69 \to 5.64$~ms for $|V| = 1{,}000 \to 5{,}000 \to 10{,}000$), consistent with the $O(|V|\log|V|)$ theoretical bound for sparse graphs. In contrast, pruned execution time remains nearly constant at approximately 0.05~ms because the threshold eliminates most of the graph regardless of size. Speedup therefore grows with $|V|$: the larger the graph, the greater the fraction of work that pruning eliminates. Peak memory for the pruned algorithm is essentially constant at approximately 4.5~KB, whereas baseline memory grows proportionally to $|V|$.

\begin{figure}[ht]
\centering
\includegraphics[width=\textwidth]{plot5_memory_scaling.png}
\caption[Peak memory scaling by graph size]{Peak memory (KB) vs.\ graph size $|V|$ at $\tau=0.01$ for baseline and pruned Dijkstra.}
\label{fig:memory_scaling}
\end{figure}

Figure~\ref{fig:memory_scaling} displays peak memory consumption on log--log axes for Erd\H{o}s--R\'{e}nyi (left) and layered (right) graphs at $\tau=0.01$. Solid lines with circle markers represent baseline Dijkstra; dashed lines with triangle markers represent the pruned variant. Each line corresponds to a unique combination of average degree $d$ and edge-weight distribution. Baseline memory scales linearly with $|V|$---from approximately 84~KB at $|V|=1{,}000$ to approximately 662~KB at $|V|=10{,}000$---consistent with the $O(|V|+|E|)$ theoretical bound. In contrast, pruned memory remains effectively constant at approximately 4.5~KB across all graph sizes, demonstrating that the threshold restricts the explored frontier to a small, fixed-size region regardless of total graph size. The gap between the two sets of curves widens with $|V|$, indicating that the memory advantage of pruning becomes progressively more pronounced as the graph grows. This constant-memory behaviour makes the pruned algorithm particularly attractive for deployment on memory-constrained systems.

\subsection{Speedup by Graph Type}

\begin{table}[h!]
\centering
\renewcommand{\arraystretch}{1.3}
\begin{tabular}{@{}lrr@{}}
\toprule
\textbf{Graph Type} & \textbf{Edge Speedup} & \textbf{Wall Speedup} \\
\midrule
Erd\H{o}s--R\'{e}nyi & 108.2$\times$ & 41.8$\times$ \\
Layered              &  57.1$\times$ & 22.5$\times$ \\
\bottomrule
\end{tabular}
\caption[Speedup by graph type]{Median speedups at $\tau = 0.01$ by graph type. ER graphs benefit more from pruning than layered graphs.}
\label{table:speedup_graph_type}
\end{table}

Erd\H{o}s--R\'{e}nyi graphs show higher speedups (41.8$\times$ wall vs.\ 22.5$\times$ for layered). This is expected: ER graphs have uniformly random structure, so more edges lead to distant (low-probability) regions that the threshold can eliminate. Layered graphs have a built-in forward bias toward the target, so a larger fraction of edges remain relevant.

\subsection{Speedup by Probability Distribution}

\begin{table}[h!]
\centering
\renewcommand{\arraystretch}{1.3}
\begin{tabular}{@{}lrr@{}}
\toprule
\textbf{Distribution} & \textbf{Edge Speedup} & \textbf{Wall Speedup} \\
\midrule
Power-law ($\alpha=2$) & 69.7$\times$ & 31.4$\times$ \\
Uniform                & 84.3$\times$ & 29.2$\times$ \\
\bottomrule
\end{tabular}
\caption[Speedup by probability distribution]{Median speedups at $\tau = 0.01$ by distribution. Both distributions achieve similar wall-clock speedups.}
\label{table:speedup_dist}
\end{table}

Both distributions achieve comparable wall-clock speedups ($\sim$30$\times$). Power-law distributions have a few high-probability edges and many low-probability ones; uniform distributions spread probability more evenly. The pruning mechanism is effective in both cases.

\subsection{Optimality Gap Analysis}

A critical finding is the \textbf{zero optimality gap} across all experiments:

\[
\Delta = \frac{|C^*_{\text{baseline}} - C^*_{\text{pruned}}|}{C^*_{\text{baseline}}} \times 100\% = 0.0000\% \quad \forall\, \text{path-found rows}
\]

This confirms Theorem~\ref{thm:all_or_nothing} (All-or-Nothing): whenever the pruned algorithm returns a path, it is the globally optimal path. The cost-based gap formula (not probability-based) eliminates floating-point rounding noise involved in $\exp(-C^*)$ conversion. Using probability-based differences would have introduced \texttt{1e-15}-scale noise that has no algorithmic significance.

\section{Real-Data Experiment Results}\label{sec:real_results}

\subsection{Datasets and Graph Statistics}

Table~\ref{table:real_data_graphs} summarises the three real-data configurations (two granularities for RetailRocket, one for RecSys~2015).

\begin{table}[h!]
\centering
\renewcommand{\arraystretch}{1.3}
\begin{tabular}{@{}lrrrl@{}}
\toprule
\textbf{Configuration} & $|V|$ & $|E|$ & \textbf{Pairs Tested} & \textbf{Source} \\
\midrule
RetailRocket (event-level)  &          3 &          9 & 20 & Kaggle~\citep{retailrocket2015} \\
RetailRocket (item-level)   &     44,711 &    101,528 & 20 & Kaggle~\citep{retailrocket2015} \\
RecSys 2015 (YOOCHOOSE)     &     12,935 &     70,442 & 20 & Kaggle~\citep{recsys2015}       \\
\bottomrule
\end{tabular}
\caption[Real-data graph configurations]{Graph sizes for the three real-data configurations. 20 random source--target pairs were tested per dataset at pruning thresholds $\tau \in \{0.0001, 0.001, 0.01, 0.1, 0.5\}$ (total: 300 rows).}
\label{table:real_data_graphs}
\end{table}

\subsection{Real-Data Speedup Results}

\begin{table}[h!]
\centering
\renewcommand{\arraystretch}{1.3}
\begin{tabular}{@{}lrrrr@{}}
\toprule
\textbf{Dataset} & \textbf{Median Speedup} & \textbf{Max Speedup} & \textbf{Paths Found} & $\Delta$ \\
\midrule
RetailRocket (event-level)  &     1.1$\times$ &     4.9$\times$ & 78 / 100 & 0.000\% \\
RetailRocket (item-level)   &   906.2$\times$ &  27,693$\times$ &  1 / 100 & 0.000\% \\
RecSys 2015 (YOOCHOOSE)     &   407.6$\times$ &   4,191$\times$ &  0 / 100 & --- \\
\bottomrule
\end{tabular}
\caption[Real-data speedup results]{Speedup results on real-world datasets. Wall-clock speedup = baseline time / pruned time. $\Delta$ is the optimality gap (0.000\% for all found paths).}
\label{table:real_data_speedup}
\end{table}

For the \textbf{RetailRocket event-level} configuration, the 3-node funnel graph (view $\to$ addtocart $\to$ transaction) is sufficiently small that both algorithms complete in microseconds; 78 of 100 configurations found paths and speedup is minimal ($\leq 5\times$) because there is almost nothing to prune. The \textbf{RetailRocket item-level} graph (44,711 nodes) yields substantial speedups up to $27{,}693\times$, although only 1 of 100 pruned configurations found a path because item-level transitions are extremely sparse and most source--target pairs become unreachable once low-probability edges are pruned. For \textbf{RecSys~2015 (YOOCHOOSE)}, the 12,935-node graph achieves a median speedup of $407.6\times$ and a maximum of $4{,}191\times$; the baseline found paths for all 100 configurations, but pruning was too aggressive at every tested $\tau$ value to preserve any complete $s\to t$ path. This reflects the sparse session structure of the YOOCHOOSE clickstream data.

\subsection{Reachability Analysis}

The low path-found rates in large real-world graphs highlight an important practical consideration. Table~\ref{table:reachability} summarises reachability across configurations.

\begin{table}[h!]
\centering
\renewcommand{\arraystretch}{1.3}
\begin{tabular}{@{}lrrl@{}}
\toprule
\textbf{Configuration} & \textbf{Baseline Found} & \textbf{Pruned Found} & \textbf{Interpretation} \\
\midrule
Synthetic (all $\tau$)       & 360 / 360 & 171 / 1,800 & Dense synthetic graphs highly connected \\
RR event-level              & 100 / 100 &  78 / 100   & 3-node funnel is fully connected \\
RR item-level               & 100 / 100 &   1 / 100   & Baseline reachable; pruning too aggressive \\
RecSys 2015                  & 100 / 100 &   0 / 100   & Baseline reachable; all pruned paths lost \\
\bottomrule
\end{tabular}
\caption[Reachability summary]{Path-found rates across synthetic and real-data experiments. The BFS connectivity guarantee in synthetic generators ensures baseline always finds paths; real-world graphs have natural fragmentation.}
\label{table:reachability}
\end{table}

In synthetic graphs, the BFS connectivity guarantee ensures the baseline always finds a path from $s$ to $t$. For the real-world datasets, the pair-selection procedure filters source--target pairs by baseline reachability, so the baseline always finds a path. The low pruned-found rates on item-level and session-level graphs reflect the aggressiveness of the $\tau$ threshold: the optimal path probability falls below $\tau$, causing all routes to be pruned even though the baseline can reach the target.

\section{Summary of Findings}

The experimental evaluation yields several principal findings. Wall-clock speedups range from $3.6\times$ at $\tau=0.001$ to $881\times$ at $\tau=0.5$ on synthetic data, and reach up to $27{,}693\times$ on real-world item-level graphs. The optimality gap remained at exactly $\Delta = 0.0000\%$ across all 250 experiment rows in which both algorithms found a path (171 of 1,800 pruned synthetic rows; 79 of 300 real-data rows), confirming the all-or-nothing property established in Theorem~\ref{thm:all_or_nothing}. Wilcoxon signed-rank tests (paired by run) confirm that all 180 synthetic configurations exhibit statistically significant speedup ($p < 0.05$); no configuration failed the significance test.

Speedup grows with graph size---from $11\times$ at $|V|=1{,}000$ to $89\times$ at $|V|=10{,}000$ at fixed $\tau=0.01$---because larger graphs contain more pruneable edges. At the same time, higher $\tau$ values reduce path-found rates from 28.6\% at $\tau=0.001$ to 3.9\% at $\tau=0.5$ on synthetic data, demonstrating the inherent reachability--speed trade-off. Real-world graphs showed further reduced pruned-found rates due to aggressive thresholding on sparse session-based structures, yet the algorithm maintained zero optimality gap on all found paths and achieved substantial speedups on large graphs.
